\section{Reliability Model: SKC Mismatch}

The taxonomy tells us \textit{how} agents fail; the reliability model tells
us \textit{why} and \textit{how likely}. At its core, \ARES{} measures the
gap between what a security task demands and what an agent can deliver across
five dimensions. The model extends classical dependability
theory~\cite{avizienis2004} by introducing knowledge incompleteness,
probabilistic reasoning errors, and capability constraints as first-class
failure sources: dimensions that do not exist in traditional deterministic
software systems.
Figure~\ref{fig:skc_pipeline_mag} summarises the end-to-end flow.

\begin{figure}[t]
\centering
\begin{tikzpicture}[
  node distance=0.35cm and 0.4cm, font=\footnotesize,
  prof/.style={rectangle, draw=black!70, rounded corners=2pt,
               fill=blue!8, minimum width=2.6cm, minimum height=0.7cm, align=center},
  mm/.style={rectangle, draw=black!70, rounded corners=2pt,
             fill=orange!12, minimum width=3.6cm, minimum height=0.7cm, align=center},
  outbox/.style={rectangle, draw=black!70, rounded corners=2pt,
              fill=green!10, minimum width=3.6cm, minimum height=0.7cm, align=center},
  arr/.style={-{Latex[length=1.8mm]}, thick, gray!70},
]
  \node[prof] (agent) {Agent profile $(\mathcal{S},\mathcal{K},\mathcal{C},\mathcal{R},\mathcal{T}_{\mathrm{age}})$};
  \node[prof, right=of agent] (task) {Task profile $(\mathcal{S}^*,\mathcal{K}^*,\mathcal{C}^*,y^*)$};
  \node[mm, below=0.4cm of $(agent.south)!0.5!(task.south)$] (delta)
       {Mismatch vector $\boldsymbol{\delta}=[\delta_S,\delta_K,\delta_C,\delta_R,\delta_T]$};
  \node[mm, below left=0.4cm and 0.05cm of delta.south] (ars)
       {ARS $= 1 - \mathbf{w}_\tau^\top\boldsymbol{\delta}$};
  \node[mm, below right=0.4cm and 0.05cm of delta.south] (eg)
       {Epistemic Gap $\mathcal{E} = \delta_K(1-U)$};
  \node[outbox, below=0.4cm of $(ars.south)!0.5!(eg.south)$] (phi)
       {Anomaly: ARS $<\theta$ \;or\; $\mathcal{E}>\epsilon$};
  \draw[arr] (agent.south) -- (delta.north);
  \draw[arr] (task.south)  -- (delta.north);
  \draw[arr] (delta.south) -- (ars.north);
  \draw[arr] (delta.south) -- (eg.north);
  \draw[arr] (ars.south)   -- (phi.north);
  \draw[arr] (eg.south)    -- (phi.north);
\end{tikzpicture}
\caption{The \ARES{} reliability pipeline. The agent's declared
profile and the task's requirement profile produce a five-dimensional
mismatch vector that feeds two parallel indicators: the Agent
Reliability Score and the Epistemic Gap. Either crossing its learned
threshold flags the execution as anomalous.}
\label{fig:skc_pipeline_mag}
\end{figure}

\subsection{Five Mismatch Dimensions}

For each dimension, we compute a normalised mismatch score between 0 (full
coverage) and 1 (total mismatch); we summarise each here and specify how it
is measured on a real agent in Section~\ref{sec:aresbench}. Full formal definitions of the
agent and task profiles and of each indicator (as set-cardinality ratios
over the task's required skills, ATT\&CK TTP subgraph, and tool manifest,
with an exponential-decay temporal term) are given in the companion
technical paper.

\begin{itemize}[noitemsep, leftmargin=*]
  \item \textbf{Skill mismatch} ($\delta_S$): the fraction of required
    procedural skills the agent lacks. For example, a log-correlation agent
    assigned binary malware analysis has a high $\delta_S$ because the
    required analytical procedure is outside its repertoire.
  \item \textbf{Knowledge mismatch} ($\delta_K$): the fraction of required
    threat knowledge missing from the agent's context, measured over the
    MITRE ATT\&CK~\cite{strom2022} knowledge graph. An agent unaware of a
    recently emerged ransomware variant has high $\delta_K$ for tasks
    involving that family.
  \item \textbf{Capability mismatch} ($\delta_C$): the fraction of
    required tools or APIs the agent cannot access. An agent that recommends
    host isolation but lacks EDR API access has high $\delta_C$: it knows
    \textit{what} to do but cannot \textit{do} it.
  \item \textbf{Reasoning mismatch} ($\delta_R$): the divergence between
    the agent's output and ground truth, capturing inference quality. This
    dimension is unique in that it can only be measured post-execution.
  \item \textbf{Temporal drift} ($\delta_T$): knowledge staleness modelled
    as exponential decay, parameterised by the threat domain's volatility.
    Ransomware intelligence ($\lambda = 0.005$~day$^{-1}$) decays rapidly;
    APT attribution knowledge ($\lambda = 0.001$~day$^{-1}$) remains stable
    for longer.
\end{itemize}

A critical design property: the first three dimensions ($\delta_S$,
$\delta_K$, $\delta_C$) and temporal drift ($\delta_T$) can all be computed
\textit{before} task execution from the agent's declared profile and the
task's requirements. In deployment, \ARES{} therefore uses skill,
knowledge, capability, and temporal drift for pre-execution
screening; reasoning mismatch ($\delta_R$) is used after execution
for diagnosis and model improvement.

\subsection{Agent Reliability Score}

The five mismatch scores are combined into a single
\textbf{Agent Reliability Score (ARS)}:
\[
  \mathrm{ARS} = 1 - \mathbf{w}_\tau^\top \boldsymbol{\delta}
\]
where $\boldsymbol{\delta} = [\delta_S, \delta_K, \delta_C, \delta_R,
\delta_T]$ is the mismatch vector and $\mathbf{w}_\tau$ is a weight vector
specific to the threat class $\tau$, learned via ridge regression on labelled
training data. An ARS of 1.0 indicates perfect alignment between agent
capability and task demand; values below a learned threshold $\theta^*$
flag the agent as unreliable for that task.

The weights are learned per threat class because different security tasks
stress different dimensions. In our experiments, knowledge coverage
($w_K = 0.35$) dominates for ransomware detection, while skill and
capability access ($w_S = 0.32$, $w_C = 0.28$) lead for DDoS mitigation,
and knowledge combined with reasoning ($w_K = 0.38$, $w_R = 0.24$) lead
for APT intrusion investigation. A single universal weight vector would
systematically misestimate reliability by at least 12\% on certain threat
classes.

\subsection{Epistemic Gap: The ``Confident Ignorance'' Problem}

The most dangerous agent failures occur not when an agent lacks knowledge,
but when it \textit{does not know that it lacks knowledge}, a failure of
confidence calibration~\cite{guo2017calibration}. We formalise this as the
\textbf{Epistemic Gap}:
\[
  \mathcal{E} = \delta_K \times (1 - U)
\]
where $U$ is the agent's expressed uncertainty (complement of its stated
confidence). An agent with a large knowledge deficit ($\delta_K \approx 1$)
that simultaneously reports high confidence ($U \approx 0$) produces a
maximal Epistemic Gap: the precondition for hallucination. Conversely, an
agent that correctly identifies its knowledge limits ($U \approx 1$) produces
a near-zero Epistemic Gap: it may fail the task but will not fabricate
an answer. Whether the dangerous regime actually arises is an empirical
question, and one much of the hallucination literature answers by
assumption. We test it directly (Empirical Insights) and find that real
frontier agents keep $U$ high as $\delta_K$ rises, so the maximal-gap
precondition is the exception rather than the rule, which makes
$\mathcal{E}$ valuable precisely as a detector of that exception.

This concept relates directly to the confidence calibration literature, which
has shown that modern neural networks (including LLMs) are systematically
overconfident on out-of-distribution inputs~\cite{guo2017calibration}. The
Epistemic Gap operationalises this finding for the specific case of agentic
AI: it measures not just whether the model is miscalibrated in general, but
whether that miscalibration occurs precisely where the agent's knowledge is
weakest: the most dangerous combination.

The distinction is critical for SOC operations. An agent that says
``I don't have sufficient context to classify this threat'' is safe: it
triggers human escalation. An agent that confidently asserts a false threat
attribution is dangerous: it triggers automated actions based on fabricated
intelligence.

\subsection{Cascade Reliability in Multi-Agent Systems}

When agents are composed into pipelines, failures propagate through
information dependencies. We model this using a directed acyclic graph
where each edge represents an inter-agent data flow, and show that cascade
reliability is bounded above by the \textit{weakest agent} in the pipeline.

This result has a counterintuitive practical consequence: decomposing
a security task into a multi-agent pipeline cannot improve reliability
beyond the weakest component agent. It can only reduce it. A triage agent
that hallucinates an IOC passes corrupted input downstream, where
enrichment inherits and amplifies it and the recommendation agent acts
on the compounded error. Our analytical results decompose this effect
cleanly. Under a matched-bucket ablation on the analytical corpus
(reported in the Empirical Insights section),
the reliability loss is attributable almost entirely to pipeline depth
rather than to the choice of agent framework, and is roughly three
times the size the naive headline gap between single-agent and
three-stage pipelines would suggest.

\subsection{From Score to Action}

Operationally, an ARS assessment runs at task-assignment time and
yields three control points. If ARS clears the threshold, the agent
runs autonomously. If it falls short, the decomposition itself
identifies the cheapest mitigation: a high $\delta_K$ calls for a
knowledge refresh, a high $\delta_C$ for tool provisioning, and a
high $\delta_T$ for a scheduled retraining cadence. The Epistemic
Gap is computed at output time and gates downstream propagation:
high-$\mathcal{E}$ outputs go to human review rather than to
automated containment or shared threat feeds. Because four of the
five dimensions are observable before execution, the screening cost
is effectively zero.
